\section{タスクの設定とデータセットの要件} \label{sec3}
本章では，家庭環境の理解を必要とする家庭内タスクについて説明し，その評価に必要なデータセットの概要を述べる．家庭環境知識とは，家庭環境内の物体に関する詳細な情報を指し，物体の現在の状態，識別番号，および物体間の関係を含む．これらの家庭環境知識に依存する家庭内タスクを評価するために必要なデータセットの要件を以下に示す．

\begin{itemize}
    \item 物体の識別番号の特定が必要なタスクとすること．これは，同じ種類の物体が複数存在する場合に，どの物体を対象とすべきかを明確にするためである．
    \item 物体の初期状態を常識的に推論すると失敗するタスクを含めること．例えば，``Turn on the light''という指示を受けた場合，常識的にはライトの状態が``OFF''であると考えられる．しかし，実際には``ON''である可能性もある．本研究では，常識知識のみに依存せず，家庭環境知識を活用することを目的とするため，このようなタスクを含める．
    \item タスクの達成条件（ゴール）が一意に定まるタスクとすること．家庭内タスクの例として``Clean up''が挙げられる．しかし，このようなタスクはゴールの定義が曖昧であり，解釈が複数存在するため，自動評価が困難である．通常は人手による評価が必要となるが，これには高い評価コストが伴うため，ゴールが明確に定義されたタスクとする．
    \item 家庭環境知識が正しく取得可能であることを前提とし，視覚情報による環境認識は行わないこと．これは，本研究において，家庭環境知識が正しく取得されたものと仮定し，それをどのようにLLMに認識させ，行動計画を行うかに焦点を当てているためである．
\end{itemize}

既存の一般的な家庭内タスクデータセットであるVHデータセットとALFREDは，上記の要件を満たしていない．まず，これらのデータセットのタスクは，特定の物体をその固有識別子によって特定することを要求していない．次に，物体の初期状態を常識的に推測するだけでは解決できないタスクが含まれていない．さらに，VHデータセットには明確に定義されたタスクの目標がないため，自動的に評価することが困難である．最後に，ALFREDデータセットは環境情報を取得するために視覚データの処理に依存しており，本研究の目的と一致していない．

そこで，本研究では，評価に必要な要件を満たす2種類のタスクを設定した．一つ目は，複数の物体を特定の状態に変化させる「状態変化タスク」であり，主に物体の状態を認識することが課題となるタスクである．もう一つは，複数の物体を特定の場所（物体の上または中）に配置する「配置タスク」であり，主に物体間の関係を認識することが課題となるタスクである．